\section{Threat Model and Goals}
\label{sec:threat}

\paragraph{Adversary.}
The adversary controls a fraction $\alpha$ of providers. It can create any
number of requester and provider identities at negligible cost (a key pair
costs about $16\,\mu$s on one core, \S\ref{sec:eval-discovery}). It can
collude across all of them, sign anything with keys it holds, publish
arbitrary receipts to any shared store, refuse to countersign, and change
behaviour over time. It controls a fraction $f$ of DHT routers and can
answer lookups arbitrarily. It cannot forge signatures of honest keys or
break SHA-256. We do not model network-level denial of service or
compromised honest keys. Key rotation and revocation are out of scope in
this version.

\paragraph{Attacks.}
We evaluate the classic reputation
attacks~\cite{hoffman2009survey,douceur2002sybil,feldman2006whitewashing},
instantiated for agent delegation:
\begin{itemize}
\item \emph{farming}: Sybil requesters issue positive receipts for their
  own providers;
\item \emph{camouflage}: farming Sybils also run truthful tasks against
  honest providers, to look like good raters;
\item \emph{collusion}: adversarial providers vouch for each other;
\item \emph{whitewashing}: adversarial providers re-enter under fresh
  identities;
\item \emph{on-off}: adversarial providers alternate good and bad
  phases;
\item \emph{slander}: Sybils file negative receipts against honest
  providers;
\item \emph{combined}: farming plus slander plus selective
  countersigning.
\end{itemize}
In discovery, the adversary attempts \emph{shard flooding} (publishing
records outside the shard they belong to), \emph{identifier grinding}
(generating identities until they land in a target shard), and
\emph{eclipse} of lookups through malicious routers.

\paragraph{Goals.}
\begin{description}
\item[G1 Attributability.] Positive evidence about a provider exists only
  if the requester agreed to it. Negative evidence is attributable to a
  requester.
\item[G2 Local soundness.] An observer's view of a provider depends only on
  receipts it can verify, weighted by its own policy. Evidence from
  issuers with no relationship to the observer contributes nothing.
\item[G3 Safe resumption.] A task resumes only from a state that the
  requester verified. A failed provider is owed at most the usage of its
  last countersigned checkpoint.
\item[G4 Bounded exposure.] Under voucher-per-checkpoint metering, neither
  side is exposed beyond one checkpoint interval of usage.
\item[G5 Composability.] Plain A2A clients keep working, and \sys{} adds
  no payment method, transport or task protocol of its own.
\end{description}
Confidentiality of task content is \emph{not} a goal. \sys{}'s
data-sensitivity zones are dual-signed labels with a refusal rule, which
provides accountability rather than confidentiality (\S\ref{sec:zones}).
